\section{礼仪的高处与粮道的低处}
\label{sec:wu-zhou-ritual-fiscal-urban-order}

高宗、武后和玄宗两度封禅泰山，武周又以改号、受命请愿和新礼仪展示权力；这些仪式确实把皇帝、官员和山川秩序连接起来，却不能被当成帝国各地同样安定的测量表。请愿文、祥瑞与正史礼志首先是政治语言，说明朝廷希望怎样被看见；边地、州县和家庭是否以同样方式理解它们，需要由碑志、文书和地方材料另行追问。%
\footnote{高宗封禅、武周建国、玄宗封禅见 ST-666-TAISHAN、ST-690-ZHOU-FOUNDATION、ST-725-TAISHAN，依据《旧唐书》《新唐书》本纪、礼仪志和《资治通鉴》。符瑞、请愿与参加规模有强烈的文体和合法性目的。}

开元九年括户检田、二十一年整顿江淮漕运，提供了从礼仪高处回到国家日常的一条路。逃户和隐漏户被重新登记，江淮粮食经转运节点送往关中，显示“盛世”必须不断修补人口、土地、仓储与道路之间的裂缝。户数和粮额在志书中往往整齐，真实的口径、地区和时点却不一致；将一组户口数字当作所有家庭生活改善的证据，正会遮住那些被重新编入赋役的人。%
\footnote{宇文融括户与裴耀卿漕运见 ST-721-YUWEN-RONG、ST-733-PEI-YAOQING，依据《旧唐书》食货志及宇文融、裴耀卿传、《新唐书》食货志和《通典》。行政数字须核对统计机构、地域与是否包含逃户。}

736年李林甫拜相、737年废杀三王、742年改州为郡，说明开元末的制度重整与继承风险同步发展。改名可以重新排列官衔与地方名目，不能消除军队供给、相权集中和皇室信任的压力。由此看，天宝以前不是一张无裂纹的“盛世”底图；它是礼仪、名册、漕运和人事暂时能够彼此配合的时期，而这套配合本身已经需要高成本维持。%
\footnote{李林甫、三王废立与天宝改制见 ST-736-LI-LINFU、ST-737-CROWN-PRINCES、ST-742-TIANBAO，依据两《唐书》玄宗纪、李林甫传、地理志与《资治通鉴》。人物归责和地方执行不能从后来的安史结局倒推。}
